% Standalone resolution manuscript generated from solution.md.
% Compile with XeLaTeX; author and review metadata are maintained in Markdown.
\documentclass[10pt,a4paper]{article}
\usepackage[margin=24mm,headheight=15pt]{geometry}
\usepackage{fontspec}
\setmainfont{DejaVuSerif.ttf}[BoldFont=DejaVuSerif-Bold.ttf,ItalicFont=DejaVuSerif-Italic.ttf,BoldItalicFont=DejaVuSerif-BoldItalic.ttf]
\setsansfont{DejaVuSans.ttf}[BoldFont=DejaVuSans-Bold.ttf,ItalicFont=DejaVuSans-Oblique.ttf]
\setmonofont{DejaVuSansMono.ttf}
\usepackage{amsmath,amssymb,unicode-math}
\setmathfont{latinmodern-math.otf}
\usepackage{xcolor,microtype,parskip,needspace,fancyhdr}
\usepackage{longtable,booktabs,array,calc}
\newcounter{none} % Pandoc's unnumbered tables
\usepackage{bookmark}
\usepackage{xurl}
\hypersetup{colorlinks=true,urlcolor=blue,linkcolor=blue,pdftitle={MI-18:
a counterexample to q-permanent monotonicity},pdfauthor={Kenta Kitamura
(formalization and supplied certificate); exposition prepared by
Codex},pdflang={en-GB}}
\urlstyle{same}
\setlength{\emergencystretch}{3em}
\providecommand{\tightlist}{\setlength{\itemsep}{0pt}\setlength{\parskip}{0pt}}
\providecommand{\pandocbounded}[1]{#1}
\setcounter{secnumdepth}{0}
\allowdisplaybreaks[1]
\widowpenalty=10000
\clubpenalty=10000
\pagestyle{fancy}
\fancyhf{}
\fancyhead[L]{\small\sffamily Counterexample manuscript / MI-18}
\fancyhead[R]{\small\sffamily 6 October 2026}
\fancyfoot[L]{\parbox[b]{0.9\textwidth}{\fontsize{8}{10}\selectfont AI-prepared
exposition of the supplied certificate; no human authorship, approval,
or discovery priority is asserted.}}
\fancyfoot[R]{\footnotesize\thepage}
\begin{document}
{\raggedright\Large\bfseries MI-18: a counterexample to q-permanent
monotonicity\par}
\vspace{0.4em}
{\normalsize\bfseries Kenta Kitamura (formalization and supplied
certificate); exposition prepared by Codex\par}
{\small \par}
\vspace{0.6em}
The universal monotonicity assertion in MI-18 is false. An ordered
family of 144 vectors in \(\mathbb C^2\) gives a positive semidefinite
Gram matrix whose \(q\)-permanent has a strictly negative derivative at
\(q=1\). Adding a sufficiently small positive multiple of the identity
gives a positive definite counterexample as well. The proof reduces the
sign of that derivative to a finite calculation with Gaussian integers;
all input integers are printed in the appendix.

The certificate and Lean formalization were supplied by Kenta Kitamura
in
\href{https://github.com/ajt60gaibb/OpenProblemsInNLA/issues/329}{issue
\#329}. The source used here is fixed at
\href{https://github.com/KitaKen1/bapat-lal-q-permanent-lean/tree/4200da4fc1a132d69c23fb877795b5b72089544b}{commit
4200da4fc1a132d69c23fb877795b5b72089544b}. Codex prepared this
conventional mathematical exposition and the accompanying independent
Python checker. This is AI-prepared prose, not a manuscript written or
approved by Kitamura. The source repository discloses AI assistance in
its formalization and certificate development; no claim of external
human review or discovery priority is made here.

\subsection{1. Statement and
construction}\label{statement-and-construction}

We index rows and columns by \(0,\ldots,n-1\). This simply shifts the
indices in MI-18 by one and preserves their order. For
\(\sigma\in S_n\), let

\nopagebreak[4]
\[
\operatorname{inv}(\sigma)
=\#\{(i,j):0\le i<j<n,\ \sigma(i)>\sigma(j)\}.
\]

For a complex matrix \(A=(A_{ij})\), put

\nopagebreak[4]
\[
P_q(A)=\sum_{\sigma\in S_n}q^{\operatorname{inv}(\sigma)}
             \prod_{i=0}^{n-1}A_{i,\sigma(i)},\qquad 0^0=1.
\tag{1}
\]

In particular, \(P_1(A)=\operatorname{per}(A)\). If \(A\) is Hermitian,
(1) is real for real \(q\): conjugating a summand replaces \(\sigma\) by
\(\sigma^{-1}\), and
\(\operatorname{inv}(\sigma^{-1})=\operatorname{inv}(\sigma)\). For the
latter equality, an inversion \(i<j\), \(\sigma(i)>\sigma(j)\)
corresponds to the inversion \((\sigma(j),\sigma(i))\) of
\(\sigma^{-1}\). Thus (1) defines an ordinary real polynomial in this
case.

Let the ordered triples \((a_j,u_j,v_j)\) be those in the appendix, and
set

\nopagebreak[4]
\[
n=144,\qquad b_j=u_j+\mathrm{i}v_j,\qquad
V_{j,:}=(a_j,b_j),\qquad H=VV^*.
\tag{2}
\]

\textbf{Theorem.} The matrix \(H\) is Hermitian positive semidefinite,
has positive diagonal, is non-diagonal, and satisfies

\nopagebreak[4]
\[
\left.\frac{d}{dq}P_q(H)\right|_{q=1}<0.
\tag{3}
\]

Moreover, there exist \(\varepsilon>0\) and \(0<q_1<q_2=1\) for which

\nopagebreak[4]
\[
A=H+\varepsilon I_{144}>0,
\qquad P_{q_2}(A)<P_{q_1}(A).
\tag{4}
\]

Consequently both the positive semidefinite version stated in MI-18 and
the positive definite version of the Bapat--Lal conjecture fail. The
construction gives order 144; it makes no assertion that this is the
smallest possible order. The proof of (4) is existential in
\(\varepsilon\) and \(q_1\).

\subsection{2. An endpoint derivative
identity}\label{an-endpoint-derivative-identity}

For an arbitrary complex \(n\times n\) matrix with \(n\ge2\), write

\nopagebreak[4]
\[
D(A)=\left.\frac{d}{dq}P_q(A)\right|_{q=1}
=\sum_{\sigma\in S_n}\operatorname{inv}(\sigma)
                         \prod_i A_{i,\sigma(i)}.
\]

If \(I=\{i,j\}\) and \(J=\{k,\ell\}\) with \(i<j\) and \(k<\ell\), write
\(A[I,J]\) for the corresponding \(2\times2\) submatrix, with these
increasing orders. The complementary submatrix is \(A[I^c,J^c]\). The
permanent of the empty matrix is understood to be 1. Define

\nopagebreak[4]
\[
T(A)=\sum_{\substack{I,J\subseteq\{0,\ldots,n-1\}\\|I|=|J|=2}}
       \det A[I,J]\,\operatorname{per}A[I^c,J^c].
\]

\textbf{Lemma 1.} For every such matrix,

\nopagebreak[4]
\[
2D(A)=\binom n2\operatorname{per}(A)-T(A).
\tag{5}
\]

\textbf{Proof.} Fix \(i<j\) and put
\(M_\sigma=\prod_r A_{r,\sigma(r)}\). Separate the permanent into the
sums \(E_{ij}\) over permutations satisfying \(\sigma(i)<\sigma(j)\) and
\(O_{ij}\) over those satisfying \(\sigma(i)>\sigma(j)\). These two
alternatives exhaust all permutations, so

\nopagebreak[4]
\[
E_{ij}+O_{ij}=\operatorname{per}(A).
\]

For each permutation in \(E_{ij}\), put \(k=\sigma(i)<\ell=\sigma(j)\)
and pair it with the permutation that exchanges the images of \(i\) and
\(j\). Their difference is

\nopagebreak[4]
\[
M_\sigma-M_{\sigma\circ(i\ j)}
=(A_{ik}A_{j\ell}-A_{i\ell}A_{jk})
                    \prod_{r\notin\{i,j\}}A_{r,\sigma(r)}.
\]

Fixing \(k<\ell\) and summing over all bijections between the remaining
row and column sets gives exactly

\nopagebreak[4]
\[
E_{ij}-O_{ij}
=\sum_{k<\ell}\det A[\{i,j\},\{k,\ell\}]
             \operatorname{per}A[\{i,j\}^c,\{k,\ell\}^c].
\]

Now sum over \(i<j\). A given \(M_\sigma\) occurs in
\(\sum_{i<j}O_{ij}\) exactly \(\operatorname{inv}(\sigma)\) times. Hence

\nopagebreak[4]
\[
T(A)=\sum_{i<j}(E_{ij}-O_{ij})
=\binom n2\operatorname{per}(A)-2D(A),
\]

as required. No positivity assumption was used.

\newpage
\subsection{3. Permanents of two-coordinate Gram
matrices}\label{permanents-of-two-coordinate-gram-matrices}

For polynomials of degree at most \(m\), define the weighted coefficient
pairing

\nopagebreak[4]
\[
\langle f,h\rangle_m
=\sum_{k=0}^{m}k!(m-k)!\,[z^k]f\,\overline{[z^k]h},
\qquad \|f\|_m^2=\langle f,f\rangle_m.
\tag{6}
\]

Consider two families of \(m\) vectors \((\alpha_i,\beta_i)\) and
\((\gamma_j,\delta_j)\), and the matrix

\nopagebreak[4]
\[
B_{ij}=\alpha_i\overline{\gamma_j}+\beta_i\overline{\delta_j}.
\]

Set \(p_X(z)=\prod_i(\alpha_i+\beta_i z)\) and
\(p_Y(z)=\prod_j(\gamma_j+\delta_j z)\).

\textbf{Lemma 2.} One has

\nopagebreak[4]
\[
\operatorname{per}(B)=\langle p_X,p_Y\rangle_m.
\tag{7}
\]

\textbf{Proof.} In the expansion of a permanent summand, let \(R\) be
the rows in which the second coordinate is chosen, and let \(C\) be the
columns matched to these rows. If \(|R|=|C|=k\), exactly \(k!(m-k)!\)
permutations map \(R\) bijectively onto \(C\) and \(R^c\) onto \(C^c\).
Their chosen-coordinate products are identical. Summing over \(R\) gives

\nopagebreak[4]
\[
\sum_{|R|=k}\prod_{i\in R}\beta_i\prod_{i\notin R}\alpha_i
=[z^k]p_X,
\]

and the column sum is the conjugate of \([z^k]p_Y\). Summing over \(k\)
proves (7), including \(m=0\) by the empty-product convention.

Return to arbitrary ordered rows \((a_j,b_j)\) of a two-column matrix
\(V\), and put

\nopagebreak[4]
\[
\ell_j(z)=a_j+b_jz,\qquad
p(z)=\prod_{j=0}^{n-1}\ell_j(z),\qquad H=VV^*.
\]

Taking the two families equal in (7) yields

\nopagebreak[4]
\[
\operatorname{per}(H)=\|p\|_n^2.
\tag{8}
\]

For \(I=\{i,j\}\), \(i<j\), write

\nopagebreak[4]
\[
w_I=a_i b_j-a_j b_i,\qquad
p_I(z)=\prod_{r\notin I}\ell_r(z),\qquad
F(z)=\sum_{|I|=2}w_I p_I(z).
\]

The \(2\times2\) submatrix \(H[I,J]\) equals \(V[I,:]V[J,:]^*\), so

\nopagebreak[4]
\[
\det H[I,J]=w_I\overline{w_J}.
\]

The complementary submatrix is likewise a mixed Gram matrix. Applying
(7) with \(m=n-2\) gives

\nopagebreak[4]
\[
\operatorname{per}H[I^c,J^c]=\langle p_I,p_J\rangle_{n-2}.
\]

Substitution into \(T(H)\) and expansion of the pairing now give

\nopagebreak[4]
\[
\begin{aligned}
T(H)
&=\sum_{I,J}w_I\overline{w_J}\langle p_I,p_J\rangle_{n-2}\\
&=\sum_{k=0}^{n-2}k!(n-2-k)!
        \left|\sum_I w_I[z^k]p_I\right|^2
=\|F\|_{n-2}^2.
\end{aligned}
\tag{9}
\]

Together, (5), (8), and (9) prove the rank-two endpoint formula

\nopagebreak[4]
\[
2D(H)=\binom n2\|p\|_n^2-\|F\|_{n-2}^2.
\tag{10}
\]

\subsection{4. The wedge polynomial and its integer
recurrence}\label{the-wedge-polynomial-and-its-integer-recurrence}

To compute \(F\) without summing over all pairs, define

\nopagebreak[4]
\[
g(z)=\sum_{j=0}^{n-1}(n-1-2j)b_j
                         \prod_{r\ne j}\ell_r(z).
\]

\textbf{Lemma 3.} \(F=-g\). In particular, \(g\) has degree at most
\(n-2\).

\textbf{Proof.} The identity

\nopagebreak[4]
\[
a_i b_j-a_j b_i=\ell_i(z)b_j-\ell_j(z)b_i
\]

implies

\nopagebreak[4]
\[
w_{\{i,j\}}p_{\{i,j\}}(z)
=b_j\prod_{r\ne j}\ell_r(z)-b_i\prod_{r\ne i}\ell_r(z).
\]

In the sum over \(i<j\), the term \(b_j\prod_{r\ne j}\ell_r\) receives
\(j\) positive contributions from indices below \(j\), and \(n-1-j\)
negative contributions from indices above it. Its net coefficient is
therefore \(2j-n+1\). This proves \(F=-g\). Each summand defining \(F\)
has degree at most \(n-2\), proving the degree assertion without
dividing by any \(\ell_j\) or coordinate.

For the certificate (2), put

\nopagebreak[4]
\[
P=\|p\|_{144}^2,\qquad S=\|g\|_{142}^2.
\]

Then (10) becomes

\nopagebreak[4]
\[
2D(H)=10296P-S,\qquad \binom{144}{2}=10296.
\tag{11}
\]

Here is a complete finite procedure for computing the two integers in
(11). Use ascending coefficients and start with \(p^{(0)}=1\) and
\(g^{(0)}=0\). For \(j=0,\ldots,143\), successively form

\nopagebreak[4]
\[
\begin{aligned}
p^{(j+1)}(z)&=(a_j+b_jz)p^{(j)}(z),\\
g^{(j+1)}(z)&=(a_j+b_jz)g^{(j)}(z)
                    +(143-2j)b_jp^{(j)}(z).
\end{aligned}
\tag{12}
\]

Induction shows that \(p^{(j)}\) is the product of the first \(j\)
factors, and \(g^{(j)}\) is the sum with one marked factor among those
first \(j\) factors, using the fixed final weights \(143-2r\). Thus
\(p^{(144)}=p\) and \(g^{(144)}=g\). More explicitly, if
\(c_{j,k}=[z^k]p^{(j)}\) and \(d_{j,k}=[z^k]g^{(j)}\), take
\(c_{0,0}=1\), all other initial coefficients zero, and use

\nopagebreak[4]
\[
\begin{aligned}
c_{j+1,k}&=a_jc_{j,k}+b_jc_{j,k-1},\\
d_{j+1,k}&=a_jd_{j,k}+b_jd_{j,k-1}
                          +(143-2j)b_jc_{j,k}.
\end{aligned}
\]

Coefficients outside their ranges are zero. Every operation is addition
or multiplication of Gaussian integers. With \((x,y)\) representing
\(x+\mathrm{i}y\), these operations are

\nopagebreak[4]
\[
(x,y)+(u,v)=(x+u,y+v),\qquad
(x,y)(u,v)=(xu-yv,xv+yu).
\]

\Needspace{9\baselineskip}
After all 144 steps, the coefficient \(d_{144,143}\) is exactly zero, as
Lemma 3 also predicts. Take

\nopagebreak[4]
\[
\begin{aligned}
P&=\sum_{k=0}^{144}k!(144-k)!
       \big((\operatorname{Re}c_{144,k})^2+(\operatorname{Im}c_{144,k})^2\big),\\
S&=\sum_{k=0}^{142}k!(142-k)!
       \big((\operatorname{Re}d_{144,k})^2+(\operatorname{Im}d_{144,k})^2\big).
\end{aligned}
\]

Evaluation on the appendix gives the exact integer inequalities

\nopagebreak[4]
\[
P>0,\qquad 10300P<S<10301P.
\tag{13}
\]

The positivity of \(P\) also follows directly: all the \(a_j\) are
nonzero, so \([z^0]p=\prod_j a_j\ne0\). The integers \(P\) and \(S\)
have 768 and 772 decimal digits, respectively; their full decimal
expansions are unnecessary for the sign test. For orientation only,
division gives

\nopagebreak[4]
\[
\frac SP\approx10300.4591950231122,\qquad
\frac{D(H)}P\approx-2.2295975115561.
\]

These decimal approximations are not used to establish (13). The exact
lower bound alone gives

\nopagebreak[4]
\[
2D(H)=10296P-S<-4P<0,
\]

which proves (3).

The machine-readable
\href{https://github.com/ajt60gaibb/OpenProblemsInNLA/blob/main/matrix-inequalities-and-norms/MI-18/certificate.json}{certificate}
contains only the ordered triples and provenance metadata. The
separately written
\href{https://github.com/ajt60gaibb/OpenProblemsInNLA/blob/main/matrix-inequalities-and-norms/MI-18/verify_certificate.py}{Python
checker} uses arbitrary-precision integers from the Python standard
library. It checks the coefficient cancellation, (13), the nonzero
off-diagonal entry used below, and an additional recurrence for \(F\).
In this second recurrence, \(p^{(j)}\) has the meaning above and
\(F^{(j)}\) is the wedge sum for the first \(j\) rows. Adjoining row
\(j\) gives

\nopagebreak[4]
\[
F^{(j+1)}
=\ell_jF^{(j)}+j b_jp^{(j)}-\ell_j\frac{d}{dz}p^{(j)},
\qquad F^{(0)}=0.
\]

Indeed the new pairs contribute
\(b_j\sum_{r<j}a_r\prod_{s<j,s\ne r}\ell_s-a_j(p^{(j)})'\), and the
first sum equals \(j p^{(j)}-z(p^{(j)})'\). This yields the displayed
recurrence. The checker verifies that its final coefficient list is
exactly the negative of that produced for \(g\) by (12).

From the repository root, the complete numerical reproduction command is

\begin{verbatim}
python3 matrix-inequalities-and-norms/MI-18/verify_certificate.py
\end{verbatim}

It returns \texttt{PASS} together with the exact rational bounds and
hashes of the two computed integers. The finite arithmetic is
independently reproducible from the appendix and (12); no copy of the
external Lean project is needed for this calculation.

\newpage
\subsection{5. Positive definiteness and strict
decrease}\label{positive-definiteness-and-strict-decrease}

For every vector \(x\in\mathbb C^{144}\),

\nopagebreak[4]
\[
x^*Hx=x^*VV^*x=\|V^*x\|_2^2\ge0,
\]

so \(H\) is Hermitian positive semidefinite. Its diagonal entries are
\(a_j^2+u_j^2+v_j^2>0\). The first two rows give

\nopagebreak[4]
\[
\begin{aligned}
H_{0,1}
&=(-60)(-55)+(72-35\mathrm{i})(80+21\mathrm{i})\\
&=9795-1288\mathrm{i}\ne0,
\end{aligned}
\]

so \(H\) is non-diagonal. This verifies every matrix hypothesis in the
positive semidefinite formulation of MI-18.

Now put \(A_\varepsilon=H+\varepsilon I_{144}\). For every
\(\varepsilon>0\) and \(x\ne0\),

\nopagebreak[4]
\[
x^*A_\varepsilon x=\|V^*x\|_2^2+\varepsilon\|x\|_2^2>0.
\]

The off-diagonal entry computed above is unaffected by this
perturbation. Also \(D(A_\varepsilon)\) is a polynomial in
\(\varepsilon\), by the finite sum defining \(D\). It is therefore
continuous at zero. Since \(D(A_0)=D(H)<0\), there exists
\(\varepsilon_0>0\) such that

\nopagebreak[4]
\[
D(A_\varepsilon)<0\qquad(0<\varepsilon<\varepsilon_0).
\]

Choose any such \(\varepsilon\), and let \(f(q)=P_q(A_\varepsilon)\).
The function \(f\) is a real polynomial with \(f'(1)<0\). By the
definition of the derivative, for all sufficiently small \(h>0\),

\nopagebreak[4]
\[
\frac{f(1)-f(1-h)}h<0.
\]

Taking \(h<1\) gives \(q_1=1-h\in(0,1)\), \(q_2=1\), and
\(f(q_2)<f(q_1)\). This proves (4). The same argument applied directly
to \(H\) gives a strict decrease for the positive semidefinite witness.

The proof specifies neither a numerical perturbation size nor a
numerical comparison point. In particular, it does not assert a decrease
for a separately proposed fixed choice such as \(10^{70}H+I\) at
\(q=1-10^{-70}\). It also gives no counterexample restricted to real
symmetric matrices. These additional statements are unnecessary to
disprove the universal complex Hermitian conjecture.

\newpage
\subsection{Appendix. The ordered integer
certificate}\label{appendix.-the-ordered-integer-certificate}

The table specifies \((a_j,u_j,v_j)\) with \(b_j=u_j+\mathrm{i}v_j\).
Each block lists a different range of row indices; use the displayed
index \(j\) to reconstruct the order \(0,1,\ldots,143\). Simultaneously
permuting rows and columns need not preserve the inversion-weighted
polynomial, so this order is part of the mathematical data.

The triples are transcribed from Kitamura's
\href{https://github.com/KitaKen1/bapat-lal-q-permanent-lean/blob/4200da4fc1a132d69c23fb877795b5b72089544b/lean/Bapat/Certificate.lean}{immutable
certificate source}. The accompanying JSON repeats precisely this small
dataset, not the external proof repository. The source repository is
distributed under
\href{https://github.com/KitaKen1/bapat-lal-q-permanent-lean/blob/4200da4fc1a132d69c23fb877795b5b72089544b/LICENSE}{Apache-2.0};
the data attribution and license provenance are recorded in the JSON.

{\small\renewcommand{\arraystretch}{0.94}\def\LTcaptype{none} % do not increment counter
\begin{longtable}[]{@{}
  >{\raggedleft\arraybackslash}p{(\linewidth - 10\tabcolsep) * \real{0.0476}}
  >{\raggedright\arraybackslash}p{(\linewidth - 10\tabcolsep) * \real{0.2857}}
  >{\raggedleft\arraybackslash}p{(\linewidth - 10\tabcolsep) * \real{0.0476}}
  >{\raggedright\arraybackslash}p{(\linewidth - 10\tabcolsep) * \real{0.2857}}
  >{\raggedleft\arraybackslash}p{(\linewidth - 10\tabcolsep) * \real{0.0476}}
  >{\raggedright\arraybackslash}p{(\linewidth - 10\tabcolsep) * \real{0.2857}}@{}}
\toprule\noalign{}
\begin{minipage}[b]{\linewidth}\raggedleft
\(j\)
\end{minipage} & \begin{minipage}[b]{\linewidth}\raggedright
\((a_j,u_j,v_j)\)
\end{minipage} & \begin{minipage}[b]{\linewidth}\raggedleft
\(j\)
\end{minipage} & \begin{minipage}[b]{\linewidth}\raggedright
\((a_j,u_j,v_j)\)
\end{minipage} & \begin{minipage}[b]{\linewidth}\raggedleft
\(j\)
\end{minipage} & \begin{minipage}[b]{\linewidth}\raggedright
\((a_j,u_j,v_j)\)
\end{minipage} \\
\midrule\noalign{}
\endhead
\bottomrule\noalign{}
\endlastfoot
0 & \((-60,72,-35)\) & 48 & \((61,58,54)\) & 96 & \((69,-10,-71)\) \\
1 & \((-55,80,-21)\) & 49 & \((57,72,40)\) & 97 & \((95,30,-1)\) \\
2 & \((52,-85,-1)\) & 50 & \((42,88,-22)\) & 98 & \((78,7,-62)\) \\
3 & \((49,-84,21)\) & 51 & \((-53,-83,-15)\) & 99 & \((84,17,-51)\) \\
4 & \((63,-61,48)\) & 52 & \((38,59,-71)\) & 100 & \((90,26,-36)\) \\
5 & \((52,-68,52)\) & 53 & \((42,76,-50)\) & 101 & \((95,27,-18)\) \\
6 & \((68,-51,53)\) & 54 & \((40,-7,-91)\) & 102 & \((92,-7,38)\) \\
7 & \((39,-83,40)\) & 55 & \((41,34,-84)\) & 103 & \((98,19,6)\) \\
8 & \((-41,91,5)\) & 56 & \((79,5,61)\) & 104 & \((66,-37,-65)\) \\
9 & \((45,-58,68)\) & 57 & \((77,29,57)\) & 105 & \((94,16,-29)\) \\
10 & \((59,-46,66)\) & 58 & \((74,46,49)\) & 106 & \((98,17,-12)\) \\
11 & \((30,-95,7)\) & 59 & \((68,65,33)\) & 107 & \((90,8,-42)\) \\
12 & \((32,-60,73)\) & 60 & \((-55,-79,27)\) & 108 & \((99,7,14)\) \\
13 & \((48,-83,-27)\) & 61 & \((65,75,12)\) & 109 & \((97,-1,25)\) \\
14 & \((49,-32,81)\) & 62 & \((61,79,-11)\) & 110 & \((-77,20,61)\) \\
15 & \((-21,79,-57)\) & 63 & \((52,-56,-64)\) & 111 & \((100,6,-1)\) \\
16 & \((-44,80,41)\) & 64 & \((50,-29,-82)\) & 112 & \((-85,6,52)\) \\
17 & \((63,-29,72)\) & 65 & \((-55,-65,53)\) & 113 & \((87,-27,41)\) \\
18 & \((36,-19,91)\) & 66 & \((77,54,33)\) & 114 & \((66,-54,-53)\) \\
19 & \((21,-23,95)\) & 67 & \((-51,-4,86)\) & 115 & \((92,-19,33)\) \\
20 & \((14,-99,-4)\) & 68 & \((75,65,15)\) & 116 & \((99,1,-15)\) \\
21 & \((-33,80,51)\) & 69 & \((-54,-38,75)\) & 117 & \((96,-2,-28)\) \\
22 & \((-73,36,-59)\) & 70 & \((84,-11,54)\) & 118 & \((92,-9,-39)\) \\
23 & \((-48,1,-88)\) & 71 & \((71,69,-10)\) & 119 & \((84,-22,-49)\) \\
24 & \((-19,75,64)\) & 72 & \((61,49,-62)\) & 120 & \((-98,12,-14)\) \\
25 & \((31,21,93)\) & 73 & \((66,63,-41)\) & 121 & \((76,-39,-52)\) \\
26 & \((59,-6,80)\) & 74 & \((84,37,40)\) & 122 & \((99,-12,3)\) \\
27 & \((2,93,-37)\) & 75 & \((86,18,48)\) & 123 & \((99,-13,-9)\) \\
28 & \((5,84,-55)\) & 76 & \((-72,-62,31)\) & 124 & \((96,-17,-24)\) \\
29 & \((16,70,70)\) & 77 & \((60,14,-79)\) & 125 & \((95,-23,20)\) \\
30 & \((45,28,85)\) & 78 & \((81,59,-3)\) & 126 & \((91,-27,-33)\) \\
31 & \((-24,-75,-61)\) & 79 & \((84,53,13)\) & 127 & \((85,-36,-39)\) \\
32 & \((27,-43,-86)\) & 80 & \((79,57,-20)\) & 128 & \((95,-27,-16)\) \\
33 & \((34,65,68)\) & 81 & \((88,41,25)\) & 129 & \((69,-61,-39)\) \\
34 & \((70,-7,72)\) & 82 & \((90,26,36)\) & 130 & \((96,-28,-2)\) \\
35 & \((-61,-18,-77)\) & 83 & \((73,45,-51)\) & 131 & \((86,-40,32)\) \\
36 & \((45,60,67)\) & 84 & \((69,30,-66)\) & 132 & \((78,-51,-36)\) \\
37 & \((54,42,73)\) & 85 & \((89,4,46)\) & 133 & \((94,-34,8)\) \\
38 & \((20,47,-86)\) & 86 & \((91,39,13)\) & 134 & \((90,-39,22)\) \\
39 & \((28,96,6)\) & 87 & \((89,45,-2)\) & 135 & \((89,-42,-17)\) \\
40 & \((76,-20,62)\) & 88 & \((61,-22,-76)\) & 136 & \((83,-52,-22)\) \\
41 & \((25,86,-44)\) & 89 & \((82,42,-39)\) & 137 & \((88,-47,-7)\) \\
42 & \((28,10,-96)\) & 90 & \((86,43,-27)\) & 138 & \((-73,64,24)\) \\
43 & \((41,-51,-76)\) & 91 & \((90,41,-15)\) & 139 & \((87,-49,7)\) \\
44 & \((44,82,37)\) & 92 & \((71,9,-70)\) & 140 & \((84,-50,22)\) \\
45 & \((71,16,68)\) & 93 & \((80,29,-53)\) & 141 & \((76,-63,-12)\) \\
46 & \((66,39,64)\) & 94 & \((95,24,20)\) & 142 & \((82,-57,10)\) \\
47 & \((42,91,6)\) & 95 & \((94,12,31)\) & 143 & \((79,-62,-1)\) \\
\end{longtable}
}

\newpage
\subsection{Sources and scope}\label{sources-and-scope}

\begin{itemize}
\tightlist
\item
  R. B. Bapat and A. K. Lal, \emph{Inequalities for the q-permanent},
  Linear Algebra and its Applications 197-198 (1994), 397-409,
  \href{https://doi.org/10.1016/0024-3795(94)90497-9}{original paper}.
\item
  L. Mitchell, \emph{A note on Bapat's q-permanent conjecture},
  Operators and Matrices 14 (2020), 915-919,
  \href{https://files.ele-math.com/articles/oam-14-56.pdf}{primary
  paper}, for the positive semidefinite formulation and earlier special
  cases.
\item
  Kenta Kitamura,
  \href{https://github.com/ajt60gaibb/OpenProblemsInNLA/issues/329}{submission
  in issue \#329}, supplying the certificate and formalization used
  here.
\item
  Kenta Kitamura,
  \href{https://github.com/KitaKen1/bapat-lal-q-permanent-lean/blob/4200da4fc1a132d69c23fb877795b5b72089544b/lean/Bapat/Main.lean}{complete
  proof entry point} and
  \href{https://github.com/KitaKen1/bapat-lal-q-permanent-lean/blob/4200da4fc1a132d69c23fb877795b5b72089544b/README.md}{source
  README}, at the fixed commit cited above. The formal target there is
  the positive definite complex Hermitian conjecture on \([-1,1]\).
\item
  \href{https://github.com/ajt60gaibb/OpenProblemsInNLA/blob/main/matrix-inequalities-and-norms/MI-18/README.md}{MI-18's
  canonical statement}. Its broader positive semidefinite formulation is
  contradicted directly by \(H\), as shown in Sections 4 and 5.
\end{itemize}

The mathematical derivation above and the integer checker establish the
counterexample without treating a reported Lean build as a premise. The
repository's separate verification record distinguishes the source and
build evidence for the external formalization from this independently
reproduced arithmetic and exposition.
\end{document}
